\expandafter\ifx\csname CAPETINCLUDEMODE\endcsname\relax
\documentclass[11pt,a4paper]{article}
% preamble.tex — packages et commandes communes a toutes les fiches CAPET SII IE
\input{glyphtounicode}
\pdfgentounicode=1
\usepackage{cmap}
\usepackage[T1]{fontenc}
\usepackage[utf8]{inputenc}
\usepackage[french]{babel}
% babel[french] active a \begin{document} des raccourcis de ponctuation (:;!?)
% qui cassent la syntaxe to[...] de circuitikz (erreur "+ or - expected") ;
% on les desactive, ce qui ne change que l'espacement fin avant ces signes.
\AtBeginDocument{\shorthandoff{:;!?}}
\usepackage[a4paper, top=1.6cm, bottom=1.6cm, left=1.4cm, right=1.4cm]{geometry}
\setlength{\headheight}{14.5pt}
\usepackage{amsmath, amssymb, amsfonts}
\usepackage{siunitx}
\sisetup{locale=FR, output-decimal-marker={,}, inter-unit-product=\ensuremath{{}\cdot{}},
         group-minimum-digits=4}
% Unites non fournies par siunitx, declarees ici une fois pour toutes (fiches + exercices).
% Les gardes \ifdefined evitent tout conflit avec une redeclaration locale restee dans un
% fichier : ne pas retirer les gardes, ne pas redeclarer ces unites dans les fichiers.
\DeclareSIUnit{\jour}{j}
\ifdefined\var\else\DeclareSIUnit{\var}{var}\fi   % puissance reactive
\ifdefined\vad\else\DeclareSIUnit{\vad}{VAD}\fi   % puissance deformante
\ifdefined\tr\else\DeclareSIUnit{\tr}{tr}\fi      % tour (tr/min)
\usepackage{array, booktabs, multirow, tabularx}
\usepackage{xcolor}
\usepackage[most, breakable]{tcolorbox}
\usepackage{multicol}
\usepackage{enumitem}
\usepackage{mathtools}
\usepackage{lmodern}
% Refonte visuelle 2026 (direction "moderne coloré") : Lato pour le texte,
% maths sans-serif assorties via newtxsf (charge apres amsmath/amssymb).
\usepackage[default]{lato}
\usepackage{newtxsf}
\usepackage{microtype}
\usepackage{fancyhdr}
\usepackage{tikz}
\usetikzlibrary{positioning,arrows.meta}
\usepackage{circuitikz}
\usepackage{lastpage}

% ── Couleurs par domaine (palette modernisée 2026) ────────────────────
\definecolor{colElec}{RGB}{29, 78, 216}        % Électrotechnique — bleu vif
\definecolor{colElecLight}{RGB}{219, 234, 254} % fond bleu clair
\definecolor{colAuto}{RGB}{21, 128, 61}        % Automatique — vert
\definecolor{colMeca}{RGB}{194, 65, 12}        % Mécanique — orange
\definecolor{colTherm}{RGB}{185, 28, 28}       % Thermique — rouge
\definecolor{colQT}{RGB}{232, 244, 233}        % (hérité, plus utilisé par questiontype)
\definecolor{colJury}{RGB}{255, 248, 225}      % fond jury (crème)
\definecolor{colJuryFrame}{RGB}{202, 138, 4}   % bordure/bandeau jury (ambre)
\definecolor{colGrayLight}{RGB}{245, 245, 247} % fond gris neutre
\definecolor{colSeparator}{RGB}{180, 180, 180} % gris séparateur
\definecolor{colExo}{RGB}{8, 145, 178}          % Banque d'exercices — cyan (distinct des fiches)
% ── Encre et gris de la refonte ───────────────────────────────────────
\definecolor{colInk}{RGB}{23, 32, 42}          % texte courant (quasi-noir)
\definecolor{colSlate}{RGB}{83, 98, 118}       % gris ardoise (méta, en-têtes)
\definecolor{colRule}{RGB}{203, 213, 225}      % filets clairs
\definecolor{colJuryBg}{RGB}{255, 248, 225}    % fond crème jury
\definecolor{colPresence}{RGB}{83, 98, 118}    % bloc 5 (présence dans les sujets)
\AtBeginDocument{\color{colInk}}

% ── Style encadré exercice + corrigé (banque d'exercices) ─────────────
\tcbset{
  exoBox/.style={
    breakable,
    enhanced,
    colback=colExo!8!white,
    colframe=colExo,
    colbacktitle=colExo,
    coltitle=white,
    fonttitle=\bfseries\small,
    left=4pt, right=4pt, top=3pt, bottom=3pt,
    before skip=5pt, after skip=5pt,
  }
}

% ── Bascule enonce / corrige (banque d'exercices) ─────────────────────
% \CAPETCORRIGE est defini en ligne de commande (pdflatex "\def\CAPETCORRIGE{1}\input{...}")
% pour compiler la variante corriges d'un meme fichier source ; par defaut (non defini),
% la compilation produit la variante enonces. \CAPETSHOW{enonce}{corrige} choisit la branche.
%
% TROIS MODES (refonte 2026-09-20, manuel de revision) :
%   - aucun drapeau   -> imprime l'enonce            (variante enonces, INCHANGE)
%   - \CAPETCORRIGE   -> imprime le corrige          (variante corriges, INCHANGE)
%   - \CAPETMANUEL    -> imprime l'enonce ET met le corrige de cote pour l'annexe
%
% Le mode manuel stocke chaque corrige dans une macro (\csgdef), sans ecriture
% fichier ni re-tokenisation : le contenu (circuitikz, accents, %) traverse intact.
% L'annexe les rejoue ensuite dans l'ordre via \CAPETANNEXE.
\usepackage{etoolbox}

\newcounter{CAPETexo}          % numero continu d'exercice dans le manuel
\setcounter{CAPETexo}{0}
\newcounter{CAPETannexe}       % compteur de relecture en annexe

% Titre courant de la notion en cours (pose par \CAPETNOTION), sert d'intertitre annexe
\newcommand{\CAPETnotioncode}{}
\newcommand{\CAPETnotiontitre}{}

% \CAPETNOTION{code}{titre} : ouvre une notion dans le manuel (memorise le libelle,
% et marque le rang du premier exercice de la notion pour l'intertitre d'annexe).
\newcommand{\CAPETNOTION}[2]{%
  \renewcommand{\CAPETnotioncode}{#1}%
  \renewcommand{\CAPETnotiontitre}{#2}%
  \csgdef{CAPETnotionstart\the\numexpr\value{CAPETexo}+1\relax}{#1 -- #2}%
}

% ── \degree robuste ───────────────────────────────────────────────────
% Certains fichiers d'exercices utilisent \degree. Ni le preambule ni siunitx
% ne le definissent, et \providecommand local ne suffit pas : sous hyperref,
% \degree existe deja comme commande factice (\ERROR), donc \providecommand
% s'abstient et l'appel echoue a la compilation du manuel (erreur invisible
% en compilation isolee). On le (re)definit donc inconditionnellement ici,
% apres le chargement des packages, pour tout le corpus.
\AtBeginDocument{\providecommand{\degree}{}\renewcommand{\degree}{\ensuremath{^\circ}}}

% Libelle decoratif : les fichiers d'exercices utilisent aussi \CAPETSHOW dans
% leur \fancyhead et leur bloc de titre, pour afficher "(enonces)" ou "(corriges)".
% Ces appels-la ne sont PAS des exercices : en mode manuel ils ne doivent ni
% consommer un numero, ni alimenter l'annexe. On les reconnait a leur premier
% argument, strictement egal au libelle ci-dessous.
\newcommand{\CAPETLIBELLEENONCES}{(\'enonc\'es)}

\ifdefined\CAPETMANUEL
  % --- Mode manuel : enonce en place, corrige differe vers l'annexe ---
  \newcommand{\CAPETSHOW}[2]{%
    \def\CAPETargun{#1}%
    \ifx\CAPETargun\CAPETLIBELLEENONCES
      #1%            appel decoratif : rendu tel quel, hors numerotation
    \else
      \CAPETSHOWEXO{#1}{#2}%
    \fi
  }
  \newcommand{\CAPETSHOWEXO}[2]{%
    \stepcounter{CAPETexo}%
    \csgdef{CAPETcorrige\theCAPETexo}{#2}%
    \edef\CAPETtmpnum{\theCAPETexo}%
    #1%
    \par\vspace{2pt}%
    {\footnotesize\itshape\color{colExo}%
     \hyperref[capet:cor:\CAPETtmpnum]{Corrige de l'exercice \CAPETtmpnum\ $\rightarrow$ p.~\pageref{capet:cor:\CAPETtmpnum}}%
     \label{capet:eno:\CAPETtmpnum}}%
  }
  % \CAPETANNEXE : rejoue tous les corriges captures, dans l'ordre, avec intertitres
  \newcommand{\CAPETANNEXE}{%
    \setcounter{CAPETannexe}{0}%
    \loop\ifnum\value{CAPETannexe}<\value{CAPETexo}%
      \stepcounter{CAPETannexe}%
      \edef\CAPETcurnum{\theCAPETannexe}%
      \ifcsname CAPETnotionstart\CAPETcurnum\endcsname
        \par\medskip
        {\normalsize\bfseries\color{colExo}\csuse{CAPETnotionstart\CAPETcurnum}}%
        \par\nobreak\smallskip
      \fi
      \begin{tcolorbox}[exoBox, breakable,
        title={Corrige de l'exercice \CAPETcurnum\
               \hfill\normalfont\footnotesize
               \hyperref[capet:eno:\CAPETcurnum]{enonce p.~\pageref{capet:eno:\CAPETcurnum}}}]%
      \label{capet:cor:\CAPETcurnum}%
      \csuse{CAPETcorrige\CAPETcurnum}%
      \end{tcolorbox}%
    \repeat
  }
\else
  % --- Modes historiques, comportement strictement inchange ---
  \newcommand{\CAPETSHOW}[2]{\ifdefined\CAPETCORRIGE #2\else #1\fi}
  \newcommand{\CAPETANNEXE}{}
\fi

% ── Difficulté 1 à 3 (étoiles pleines, amssymb) ───────────────────────
\newcommand{\difficulte}[1]{%
  \ifcase#1\relax
  \or $\bigstar$\;niveau 1/3%
  \or $\bigstar\bigstar$\;niveau 2/3%
  \or $\bigstar\bigstar\bigstar$\;niveau 3/3%
  \fi
}

% ── Style encadré extrait de document technique (banque d'exercices) ──
\tcbset{
  docTech/.style={
    breakable,
    enhanced,
    colback=colGrayLight,
    colframe=colSeparator,
    fonttitle=\bfseries\small\itshape,
    coltitle=black,
    colbacktitle=colGrayLight,
    left=4pt, right=4pt, top=3pt, bottom=3pt,
    before skip=4pt, after skip=4pt,
  }
}

% ── Style tcolorbox de section principale ────────────────────────────
% Refonte 2026 : bandeau titre coloré arrondi + fond très pâle, sans cadre.
\tcbset{
  ficheSection/.style={
    breakable, enhanced jigsaw,
    frame hidden, boxrule=0pt, arc=1.8mm,
    colback=#1!4!white,
    colbacktitle=#1,
    coltitle=white,
    fonttitle=\sffamily\bfseries,
    left=7pt, right=7pt, top=5pt, bottom=5pt,
    toptitle=3pt, bottomtitle=3pt,
    before skip=10pt, after skip=8pt,
  }
}

% ── Style encadré question type ──────────────────────────────────────
% Carte blanche, bandeau sombre, liseré très clair.
\tcbset{
  questiontype/.style={
    breakable, enhanced jigsaw,
    colback=white,
    frame hidden, boxrule=0pt, arc=1.8mm,
    colbacktitle=colElec!78!black,
    coltitle=white,
    fonttitle=\sffamily\bfseries\small,
    left=7pt, right=7pt, top=4pt, bottom=4pt,
    toptitle=2.5pt, bottomtitle=2.5pt,
    before skip=8pt, after skip=8pt,
    borderline={0.6pt}{0pt}{colElec!12!white},
  }
}

% ── Style encadré jury ───────────────────────────────────────────────
% Bandeau ambre, fond crème, sans cadre.
\tcbset{
  jury/.style={
    breakable, enhanced jigsaw,
    colback=colJuryBg,
    frame hidden, boxrule=0pt, arc=1.8mm,
    colbacktitle=colJuryFrame,
    coltitle=white,
    fonttitle=\sffamily\bfseries,
    left=7pt, right=7pt, top=4pt, bottom=4pt,
    toptitle=3pt, bottomtitle=3pt,
    before skip=10pt, after skip=8pt,
  }
}

% ── En-tête / pied de page communs (refonte 2026) ─────────────────────
% #1 code fiche, #2 titre court, #3 couleur domaine
\newcommand{\ficheHeaders}[3]{%
  \fancyhead[L]{\small\sffamily{\bfseries\textcolor{#3}{#1}}\hspace{0.7em}\textcolor{colSlate}{#2}}%
  \fancyhead[R]{\small\sffamily\textcolor{colSlate}{CAPET SII IE\, --- \,Session 2026}}%
  \fancyfoot[C]{\small\sffamily\textcolor{colSlate}{\thepage\,/\,\pageref{LastPage}}}%
  \fancyhead[C]{}\fancyfoot[L]{}\fancyfoot[R]{}%
}

% ── Bloc en-tête de fiche (BLOC 0, refonte 2026) ──────────────────────
% #1 code  #2 titre  #3 couleur domaine  #4 domaine  #5 ligne présence  #6 sous-thèmes
\newcommand{\ficheEntete}[6]{%
  \begin{tcolorbox}[enhanced, breakable,
    frame hidden, boxrule=0pt, arc=2.5mm,
    interior style={left color=#3, right color=#3!75!black},
    left=9pt, right=9pt, top=7pt, bottom=7pt,
    before skip=0pt, after skip=6pt]
    {\color{white}\sffamily\bfseries\LARGE #1\quad #2}\\[4pt]
    {\color{#3!22!white}\sffamily\small \MakeUppercase{#4}\quad\textbar\quad #5}\\[1pt]
    {\color{#3!35!white}\sffamily\footnotesize Sous-thèmes : #6}
  \end{tcolorbox}%
}

% ── Mise en page ─────────────────────────────────────────────────────
\pagestyle{fancy}
\fancyhf{}
\renewcommand{\headrulewidth}{0.5pt}
\renewcommand{\footrulewidth}{0pt}
\AtBeginDocument{%
  \renewcommand{\headrule}{{\color{colRule}\hrule width\headwidth height 0.5pt}}%
}

\setlist[itemize]{leftmargin=1.2em, itemsep=1pt, topsep=2pt, parsep=0pt}
\setlist[enumerate]{leftmargin=1.2em, itemsep=3pt, topsep=2pt, parsep=0pt}

\setlength{\parskip}{2pt}
\setlength{\columnsep}{10pt}

\def\CAPETBEGINFILE{\begin{document}}
\def\CAPETENDFILE{\end{document}}
\else
\def\CAPETBEGINFILE{}
\def\CAPETENDFILE{}
\fi

\ficheHeaders{T01}{Grandeurs électrocinétiques}{colElec}

\CAPETBEGINFILE

% -----------------------------------------------------------------------
%  BLOC 0 - EN-TETE
% -----------------------------------------------------------------------

\ficheEntete{T01}{Grandeurs électriques fondamentales et électrocinétique}{colElec}%
  {électrotechnique}%
  {\textbf{Présence :} \textbf{13/13} sujets (2013-2026) - thème \emph{systématiquement} présent}%
  {lois de Kirchhoff - diviseur de tension/courant - valeur efficace - régime continu vs alternatif}

% -----------------------------------------------------------------------
%  BLOC 1 - RAPPELS DE COURS
% -----------------------------------------------------------------------

\begin{tcolorbox}[
  ficheSection=colElec,
  title={1. Rappels de cours},
]

% -- 1.1 Lois fondamentales ------------------------------------------
\textbf{1.1 Lois de Kirchhoff et loi d'Ohm}

\smallskip
Tout circuit peut être résolu par application systématique de ces trois lois.

\smallskip
\begin{center}
\renewcommand{\arraystretch}{1.45}
\begin{tabular}{llll}
\toprule
\textbf{Loi} & \textbf{Énoncé} & \textbf{Expression} & \textbf{Remarque} \\
\midrule
Loi d'Ohm & Tension = résistance $\times$ courant & $u = R \cdot i$ & $R$ en $\Omega$ \\[3pt]
Loi d'Ohm généralisée & Avec source de tension $e$ & $u = R \cdot i + e$ & Branche active \\[3pt]
Loi des mailles (LMK) & Somme algébrique des tensions & $\sum e_k = \sum R_k i_k$ & Choisir un sens \\[3pt]
Loi des nœuds (LNK) & Conservation des courants & $\sum i_\text{entrants} = \sum i_\text{sortants}$ & En tout nœud \\
\bottomrule
\end{tabular}
\end{center}

\smallskip
\textbf{Diviseurs :}
\[
  V_s = V_e \cdot \frac{R_B}{R_A + R_B} \quad \text{(tension)}
  \qquad\qquad
  I_1 = I \cdot \frac{R_2}{R_1 + R_2} \quad \text{(courant)}
\]

\smallskip
\begin{center}
\begin{circuitikz}[european, scale=0.85]
  \draw (0,0) to[V, l=$E$] (0,3)
        to[R, l=$R_A$] (3,3)
        to[R, l_=$R_B$] (3,0) node[circ]{}
        -- (0,0) node[circ]{};
  \draw (3,1.5) to[short] (4.3,1.5) node[right] {$V_s$};
\end{circuitikz}
\end{center}

\medskip\hrule\vspace{4pt}

% -- 1.2 Regime sinusoidal ------------------------------------------
\textbf{1.2 Grandeurs sinusoïdales - valeurs efficaces}

\smallskip
Pour $u(t) = U_\text{max}\cos(\omega t + \phi)$, la \textbf{valeur efficace} vaut :
\[
  U_\text{eff} = \frac{U_\text{max}}{\sqrt{2}} \quad \text{(sinusoïde pure)}
  \qquad
  U_\text{eff} = \sqrt{\frac{1}{T}\int_0^T u^2(t)\,dt} \quad \text{(cas général)}
\]

Réseau français : $U_\text{eff} = \SI{230}{\volt}$ (phase-neutre), $U_\text{max} = 230\sqrt{2} \approx \SI{325}{\volt}$.

\smallskip
\begin{center}
\begin{tikzpicture}[scale=0.9, font=\footnotesize]
  \draw[->] (0,-1.5) -- (0,1.6) node[above] {$u(t)$};
  \draw[->] (-0.3,0) -- (6.3,0) node[right] {$t$};
  \draw[thick, colElec, domain=0:6.28, samples=80, smooth]
    plot (\x, {1.3*cos(\x r)});
  \draw[dashed] (0,1.3) -- (6.2,1.3) node[right] {\scriptsize $U_\text{max}$};
  \draw[dashed] (0,0.92) -- (6.2,0.92) node[right] {\scriptsize $U_\text{eff}=U_\text{max}/\sqrt{2}$};
  \draw[dashed] (6.28,0) -- (6.28,-1.5);
  \node[below] at (6.28,-1.5) {$T$};
\end{tikzpicture}
\end{center}

\medskip\hrule\vspace{4pt}

% -- 1.3 ALI parfait ----------------------------------------------
\textbf{1.3 Amplificateur Linéaire Intégré (ALI) parfait - hypothèses}

\smallskip
En régime linéaire, un ALI parfait vérifie deux hypothèses fondamentales :

\begin{center}
\renewcommand{\arraystretch}{1.4}
\begin{tabular}{lll}
\toprule
\textbf{Hypothèse} & \textbf{Expression} & \textbf{Conséquence} \\
\midrule
Entrées au même potentiel & $v_+ = v_-$ & Calculer $v_-$ depuis $v_+$ \\[3pt]
Courants d'entrée nuls & $i_+ = i_- = 0$ & LNK simplifiée aux entrées \\
\bottomrule
\end{tabular}
\end{center}

\smallskip
Application : dans un montage inverseur, la LNK au nœud $v_-$ donne directement le gain $G = -R_2/R_1$.

\end{tcolorbox}

% -----------------------------------------------------------------------
%  BLOC 2 - FORMULAIRE SYNTHETIQUE
% -----------------------------------------------------------------------

\begin{tcolorbox}[
  ficheSection=colElec,
  title={2. Formulaire synthétique},
]

\small
\begin{multicols}{2}
\renewcommand{\arraystretch}{1.45}
\begin{tabular}{@{}p{4.2cm}ll@{}}
\toprule
\textbf{Formule} & \textbf{Usage} & \textbf{U.} \\
\midrule
\multicolumn{3}{@{}l}{\footnotesize\textit{Lois de base (DC et AC)}} \\
$u = R \cdot i$ & Loi d'Ohm & V \\
$u = R \cdot i + e$ & Branche active & V \\
$V_s = V_e \cdot R_B/(R_A+R_B)$ & Div. tension & V \\
$I_1 = I \cdot R_2/(R_1+R_2)$ & Div. courant & A \\
\midrule
\multicolumn{3}{@{}l}{\footnotesize\textit{Valeur efficace}} \\
$U_\text{eff} = U_\text{max}/\sqrt{2}$ & Sinusoïde & V \\
$E_C = \tfrac{1}{2}C V^2$ & Énergie condo & J \\
$E_L = \tfrac{1}{2}L I^2$ & Énergie bobine & J \\
\midrule
\multicolumn{3}{@{}l}{\footnotesize\textit{Constantes de temps}} \\
$\tau = RC$ & Circuit RC & s \\
$\tau = L/R$ & Circuit RL & s \\
\bottomrule
\end{tabular}

\columnbreak

\renewcommand{\arraystretch}{1.45}
\begin{tabular}{@{}p{4.2cm}ll@{}}
\toprule
\textbf{Formule} & \textbf{Usage} & \textbf{U.} \\
\midrule
\multicolumn{3}{@{}l}{\footnotesize\textit{Puissances}} \\
$P = U \cdot I$ & Continu & W \\
$P = U I\cos\varphi$ & AC monophasé & W \\
$P = \sqrt{3}\,UI\cos\varphi$ & AC triphasé & W \\
$Q = UI\sin\varphi$ & Réactive & VAR \\
$S = UI = \sqrt{P^2+Q^2}$ & Apparente & VA \\
\midrule
\multicolumn{3}{@{}l}{\footnotesize\textit{Moteur asynchrone (MAS)}} \\
$N_s = 60f/p$ & Vitesse synchro & tr/min \\
$g = (N_s - N)/N_s$ & Glissement & - \\
\midrule
\multicolumn{3}{@{}l}{\footnotesize\textit{Réseau FR}} \\
$V = \SI{230}{\volt}$, $U = \SI{400}{\volt}$ & Phase / Ligne & V \\
$\omega = 2\pi\times50 \approx 314$ & Pulsation & rad/s \\
\bottomrule
\end{tabular}

\end{multicols}
\normalsize

\end{tcolorbox}

% -----------------------------------------------------------------------
%  BLOC 3 - QUESTIONS TYPES
% -----------------------------------------------------------------------

\begin{tcolorbox}[
  ficheSection=colElec,
  title={3. Questions types - extraites des sujets réels},
]

% -- QT1 -------------------------------------------------------------

\begin{tcolorbox}[questiontype,
  title={QT1 - Loi des mailles avec ALI parfait \hfill\normalfont\footnotesize E2 disciplinaire IE 2024}
]

\textbf{Énoncé.}
Le montage comporte un ALI parfait en régime linéaire.
La résistance $R_1$ est connectée entre l'entrée $e_{AB}$ et l'entrée inverseuse $v_-$.
La résistance $R_2$ est en contre-réaction (de la sortie vers $v_-$).
Écrire la relation liant $e_{AB}$ à la tension aux bornes de $R_2$.

\tcbline

\textbf{Corrigé.}

\textbf{Étape 1 — Hypothèse ALI.}
En régime linéaire : $v_+ = v_- = 0$ (entrée non-inverseuse à la masse) et $i_- = 0$.

\textbf{Étape 2 — LNK au nœud $v_-$.}
Le courant entrant par $R_1$ doit ressortir par $R_2$ (puisque $i_- = 0$) :
\[
  \frac{e_{AB} - 0}{R_1} = \frac{0 - V_s}{R_2}
  \quad\Longrightarrow\quad
  V_s = -\frac{R_2}{R_1}\,e_{AB}
\]

\textbf{Conclusion.} La tension $e_{AB}$ se retrouve intégralement traduite en sortie avec le gain $-R_2/R_1$ : le montage se comporte en amplificateur inverseur linéaire, conforme au cahier des charges d'un étage de conditionnement.
\end{tcolorbox}

\vspace{4pt}

% -- QT2 -------------------------------------------------------------

\begin{tcolorbox}[questiontype,
  title={QT2 - Diviseur de tension - rapport d'atténuation \hfill\normalfont\footnotesize E2 disciplinaire IE 2023}
]

\textbf{Énoncé.}
Le circuit F2 est un pont diviseur de tension alimenté par $V_{e1}$.
Il comporte deux résistances $R_A$ et $R_B$ ($R_A = 9\,R_B$) montées en série.
Exprimer $V_{s1}$ en fonction de $V_{e1}$, $R_A$, $R_B$ puis calculer le rapport $V_{s1}/V_{e1}$.

\tcbline

\textbf{Corrigé.}

\textbf{Étape 1 — Formule du diviseur de tension.}
\[
  V_{s1} = V_{e1} \cdot \frac{R_B}{R_A + R_B}
\]

\textbf{Étape 2 — Application numérique avec $R_A = 9\,R_B$.}
\[
  V_{s1} = V_{e1} \cdot \frac{R_B}{9R_B + R_B} = V_{e1} \cdot \frac{1}{10}
  \quad\Longrightarrow\quad
  \frac{V_{s1}}{V_{e1}} = \mathbf{\frac{1}{10}}
\]

\textbf{Conclusion.} Le circuit réalise une atténuation d'un facteur 10 : il ramène une tension $V_{e1}$ à un dixième de sa valeur, ce qui correspond à l'adaptation de niveau attendue en entrée du capteur/CAN du sujet.
\end{tcolorbox}

\vspace{4pt}

% -- QT3 -------------------------------------------------------------

\begin{tcolorbox}[questiontype,
  title={QT3 - Constante de temps d'un circuit RL - modèle Thévenin \hfill\normalfont\footnotesize E2 disciplinaire IE 2023}
]

\textbf{Énoncé.}
Une batterie est modélisée par une force électromotrice $E_{Li}$ en série avec une résistance interne $R_{Li}$ et une inductance $L$.
Le circuit est fermé sur une charge résistive $R$.
Calculer la constante de temps $\tau$ du circuit en régime transitoire.

\tcbline

\textbf{Corrigé.}

\textbf{Étape 1 — Équation différentielle par LMK.}
La loi des mailles donne :
\[
  E_{Li} = R_{Li}\,i + L\,\frac{di}{dt} + R\,i = (R_{Li}+R)\,i + L\,\frac{di}{dt}
\]

\textbf{Étape 2 — Identification de la constante de temps.}
On identifie à la forme canonique du premier ordre :
\[
  \tau\,\frac{di}{dt} + i = \frac{E_{Li}}{R_{Li}+R}
  \quad\Longrightarrow\quad
  \boxed{\tau = \frac{L}{R_{Li}+R}}
\]

\textbf{Conclusion.} En négligeant la charge ($R \ll R_{Li}$), on retrouve $\tau \approx L/R_{Li}$ : le temps de réponse est piloté par la résistance interne de la batterie, un point de dimensionnement clé pour un système embarqué.
\end{tcolorbox}

\end{tcolorbox}

% -----------------------------------------------------------------------
%  BLOC 4 - ATTENTES DU JURY
% -----------------------------------------------------------------------

\begin{tcolorbox}[
  jury,
  title={4. Attentes du jury - Erreurs récurrentes et points de vigilance},
]

\begin{enumerate}[label=\textbf{\arabic*.}, leftmargin=*, itemsep=5pt]

  \item \textbf{Confondre puissance continue et alternative.}
  En DC : $P = U \cdot I$. En AC : $P = U_\text{eff} I_\text{eff} \cos\varphi$.
  Utiliser $P = UI$ en régime AC revient à supposer $\cos\varphi = 1$, ce qui est rarement vrai.

  \item \textbf{Oublier de choisir un sens de parcours avant d'écrire une LMK.}
  Une tension est comptée positivement si le sens de parcours va du $-$ au $+$ de la source.
  Ne jamais écrire une LMK sans avoir tracé le sens sur le schéma.

  \item \textbf{Confondre valeur de crête et valeur efficace.}
  $U_\text{max} = U_\text{eff}\sqrt{2}$ uniquement pour un signal sinusoïdal pur.
  Pour un signal déformé : utiliser la définition intégrale.

  \item \textbf{Oublier les hypothèses de l'ALI parfait.}
  $v_+ = v_-$ et $i_+ = i_- = 0$ en régime linéaire : ces deux hypothèses doivent être
  invoquées explicitement avant d'appliquer la LNK, sans quoi la démarche n'est pas valide.

  \item \textbf{Confondre tension de phase et tension composée en triphasé.}
  Réseau 400\,V/230\,V : $V = \SI{230}{\volt}$ (phase-neutre), $U = \SI{400}{\volt}$ (entre phases).
  $P = \sqrt{3}\,UI\cos\varphi$ utilise la tension \emph{composée} $U$ et le courant \emph{de ligne}.

  \item \textbf{Ne pas justifier les étapes de calcul.}
  Le jury attend la démarche, pas seulement le résultat numérique.
  Écrire la formule littérale, substituer les valeurs, isoler l'inconnue, encadrer le résultat.

\end{enumerate}

\medskip
\begin{tcolorbox}[
  colback=colJuryFrame!12!colJury,
  colframe=colJuryFrame,
  left=4pt, right=4pt, top=2pt, bottom=2pt,
  fontupper=\small\itshape
]
Rapport jury 2024 : «~Le jury a constaté le manque de maîtrise des formules de base sur les réseaux triphasés, le régime sinusoïdal et les relations d'électrocinétique telles que la loi des mailles, la loi des nœuds ou encore la loi d'Ohm.~»\\
Rapport jury 2025 : «~De nombreux candidats ne maîtrisent pas les lois basiques de l'électricité (la loi des nœuds, la loi des mailles, les calculs de puissance ou encore la loi d'Ohm).~»
\end{tcolorbox}

\end{tcolorbox}

% -----------------------------------------------------------------------
%  BLOC 5 - PRESENCE DANS LES SUJETS
% -----------------------------------------------------------------------

\begin{tcolorbox}[
  ficheSection=colPresence,
  title={Présence dans les sujets analysés (2013-2026)},
]
\begin{center}
\footnotesize\renewcommand{\arraystretch}{1.15}
\begin{tabular}{lccccccccccccc}
\toprule
\textbf{Année} & 13 & 14 & 15 & 16 & 17 & 18 & 19 & 20 & 21 & 22 & 23 & 24 & 25 \\
\midrule
E1 (commune)    & $\checkmark$ & $\checkmark$ & $\checkmark$ & $\checkmark$ & $\checkmark$ & $\checkmark$ & $\checkmark$ & $\checkmark$ & $\checkmark$ & $\checkmark$ & $\checkmark$ & $\checkmark$ & $\checkmark$ \\
E2 discipl. IE  & $\checkmark$ & $\checkmark$ & $\checkmark$ & $\checkmark$ & $\checkmark$ & $\checkmark$ & $\checkmark$ & $\checkmark$ & $\checkmark$ & $\checkmark$ & $\checkmark$ & $\checkmark$ & $\checkmark$ \\
\bottomrule
\end{tabular}
\end{center}
\smallskip
\footnotesize\textbf{Sous-thèmes fréquents :}
Loi des mailles / nœuds (13/13) - loi d'Ohm (13/13) - diviseur de tension (10/13) - ALI parfait (8/13) - constante de temps RC/RL (7/13) - valeur efficace (6/13)

\medskip
\textbf{Ordres de grandeur (repères rapides).}
\begin{itemize}[leftmargin=1.3em]
  \item Réseau BT FR : tension phase-neutre typique \SI{230}{\volt}, tension phase-phase typique \SI{400}{\volt}, fréquence \SI{50}{\hertz}.
  \item Courants d'étude fréquents en sujet : de quelques \si{\ampere} à quelques dizaines de \si{\ampere}.
  \item En conversion d'énergie, un rendement attendu reste généralement entre 0{,}8 et 0{,}98 selon la technologie.
\end{itemize}

\medskip
\textbf{Synthèse express (à refaire sans support).}
\begin{enumerate}[leftmargin=1.4em]
  \item Poser les hypothèses du modèle et les conventions de signe.
  \item Écrire 2 à 4 relations clés avec unités.
  \item Réaliser un calcul numérique complet et vérifier l'ordre de grandeur.
  \item Justifier le résultat par un argument technique de niveau CAPET.
\end{enumerate}

\end{tcolorbox}

\CAPETENDFILE
